% 【本文件写什么】一级组件 wateros-driver（§ 级聚合）
%   - 组件职责与在内核中的位置：块设备、字符设备、网络设备与板级驱动聚合入口。
%   - 聚合 crate os/components/wateros-driver/src/lib.rs 的 pub mod 树与对外导出面
%   - 子模块划分、与相邻组件的依赖边界（谁依赖谁、走哪条门面）
%   - 根 feature / 本组件 Cargo.toml feature 如何选用各 impl
%   - 1～2 段综述后，由下方 \input 展开子模块；不在此逐条列举 api trait
% 【事实来源】
%   - os/components/wateros-driver/src/lib.rs、Cargo.toml
%   - docs/exports/public-api/wateros-driver.md
%   - docs/exports/features/wateros-driver.md
%   - os/feature-tree.txt 中 wateros-driver 相关段
% 【不写什么】
%   - api-v0 契约细节 → 子目录 api/api-v0.tex
%   - 各 impl 算法/平台细节 → 子目录 impl/*.tex

\section{设备驱动组件}

驱动组件把板级探测、块设备、字符设备和网络设备收成内核能统一注册的入口。引导早期只记下设备树物理地址；内存子系统就绪后再扫描总线、握手 VirtIO 块卡与网卡，并把串口、RTC 等字符设备登记进全局表。探测失败只打日志，不阻断内核继续启动。

RISC-V 主线从设备树找 MMIO VirtIO 与 NS16550；LoongArch 主线走 PCIe 配置空间枚举 VirtIO-PCI，早期 UART 用固定 MMIO 窗口。块设备可选一层写穿缓存；网卡注册后由内置协议栈周期性轮询，再与用户态 socket 路径衔接。设备就绪后文件系统组件刷新设备目录，根卷探测才能看到块设备节点。当前 bring-up 仍以「第一块盘、第一张网卡」为主，多实例表已具备但未全面用到。

\begin{rustcode}
// driver 聚合：两步引导 + 平台 feature 选路
pub fn init_when_boot(dtb_pa: usize) { active_impl::init_when_boot(dtb_pa); }
pub fn init_after_boot() {
    // RISC-V：DTB → virtio-mmio blk/net + UART
    // LoongArch：PCIe ECAM → virtio-pci blk/net + 板级 UART
    // 末尾 devfs::refresh()，失败仅 warn
}
// block：BlockDevice、register_block_device、可选 CachingBlockDevice
// network::stack：init / poll / TCP-UDP socket 句柄
// character：null、RTC 等内置桩
\end{rustcode}
